\section{Conclusion}
 In this work, we proposed a definition of fairness based upon equity, demonstrated its appeal as a fairness outcome to a wide audience, and formalized it for classification. We tested this approach in a traditional cross validation setup, and demonstrated how it can be used in a real-world environment, such as unfairness that can arise from the feedback loop. Our results show the effectiveness of our method in mitigating bias and achieving fairness. We also performed human evaluation to evaluate our notion in different scenarios with the equality/parity notion of fairness. As a future direction, our definition can be utilized to achieve and study the effects of equity in classification with different techniques. In this work, we provide a framework for equity to be formalized; however, there is still work to be done in the area of fairness with regards to equity. Future work is to further study how the equity notion interacts with other existing definitions of fairness, such as equality of opportunity, equalized odds or other definitions in the equality domain other than statistical parity. It can also be extended to other machine learning tasks such as regression.

\section{Acknowledgments} 
We wanted to thank Hrayr Harutyunyan and Mozhdeh Gheini for their help and comments.

